%------------------------------------------------------------------------------%
\documentclass[fleqn,utf8,aspectratio=169,12pt]{beamer}

\usepackage{gaal}

\title{Ortogonalização de Gram--Schmidt}

%------------------------------------------------------------------------------%
\begin{document}

\addtitle

%------------------------------------------------------------------------------%
\section{Projeção ortogonal}

%------------------------------------------------------------------------------%
\begin{frame}{Projeção Ortogonal}
  A projeção de $\vec{v}$ sobre $\vec{w}$ é
  \[
    \pause
    \vec{v}_{\parallel}
    \pause
    \;=\; \proj_{\vec{w}}(\vec{v})
    \pause
    \;=\; \frac{\vec{v}\cdot\vec{w}} {\norm{\vec{w}}^{2}} \; \vec{w}
  \]

  \pause
  A componente perpendicular é
  \[
    \pause
    \vec{v}_{\perp}
    \pause
    \;=\; \vec{v} - \proj_{\vec{w}}(\vec{v})
  \]
\end{frame}

%------------------------------------------------------------------------------%
\section{Ortogonalização de Gram--Schmidt}

%------------------------------------------------------------------------------%
\begin{frame}{Gram--Schmidt}
  Dados vetores \emph{linearmente independentes}
  \[
    \pause
    \vec{v}_1, \vec{v}_2, \ldots, \vec{v}_n
  \]
  \pause
  o processo de \emph{Gram--Schmidt} constrói vetores \emph{ortogonais}
  \[
    \pause
    \vec{u}_1, \vec{u}_2, \ldots, \vec{u}_n
  \]
  \pause
  que podem ser normalizados para produzir uma \emph{base ortonormal}
  \[
    \pause
    \vec{e}_1, \vec{e}_2, \ldots, \vec{e}_n
  \]
  \pause
  com \;
  \(
    \pause
    \ger \{\vec{e}_1,\ldots,\vec{e}_k\}
    \;=\; \ger \{\vec{u}_1,\ldots,\vec{u}_k\}
    \;=\; \ger \{\vec{v}_1,\ldots,\vec{v}_k\}
  \)
\end{frame}

%------------------------------------------------------------------------------%
\begin{frame}{Gram--Schmidt: Orimeiro Vetor}
  Escolha
  \[
    \vec{u}_1 = \vec{v}_1
  \]

  \pause
  Normalize
  \[
    \vec{e}_1 = \frac{\vec{u}_1}{\norm{\vec{u}_1}}
  \]
\end{frame}

%------------------------------------------------------------------------------%
\begin{frame}{Gram--Schmidt: segundo vetor}
  \onslide<+->{Retire de $\vec{v}_2$ sua componente na direção de $\vec{u}_1$}
  \begin{align*}
    \onslide<+->{\vec{u}_2}
    \onslide<+->{& = \vec{v}_2 - \proj_{\vec{u}_1}(\vec{v}_2)}
    \\
    \onslide<+->{& = \vec{v}_2
      - \frac{\vec{v}_2\cdot\vec{u}_1} {\norm{\vec{u}_1}^{2}} \; \vec{u}_1}
  \end{align*}
  \onslide<+->{Depois normalize}
  \[
    \onslide<+->{\vec{e}_2 = \frac{\vec{u}_2}{\norm{\vec{u}_2}}}
  \]
\end{frame}

%------------------------------------------------------------------------------%
\begin{frame}{Gram--Schmidt: terceiro vetor}
  \onslide<+->{Retire de $\vec{v}_3$ as componentes nas direções anteriores}
  \begin{align*}
    \onslide<+->{\vec{u}_3}
    \onslide<+->{& = \vec{v}_3
      - \proj_{\vec{u}_1}(\vec{v}_3)
      - \proj_{\vec{u}_2}(\vec{v}_3)}
    \\
    \onslide<+->{& = \vec{v}_3
      - \frac{\vec{v}_3\cdot\vec{u}_1}{\norm{\vec{u}_1}^{2}} \; \vec{u}_1
      - \frac{\vec{v}_3\cdot\vec{u}_2}{\norm{\vec{u}_2}^{2}} \; \vec{u}_2}
  \end{align*}
  \onslide<+->{Depois normalize}
  \[
    \onslide<+->{\vec{e}_3 = \frac{\vec{u}_3}{\norm{\vec{u}_3}}}
  \]
\end{frame}

%------------------------------------------------------------------------------%
\begin{frame}{Gram--Schmidt em $\R^n$}
  \onslide<+->{Para $k\geq 2$}
  \onslide<+->{\[
    \vec{u}_k
    = \vec{v}_k
    - \sum_{j=1}^{k-1} \; \frac{\vec{v}_k\cdot\vec{u}_j} {\norm{\vec{u}_j}^{2}} \; \vec{u}_j
  \]}
  \onslide<+->{Depois
  \[
    \vec{e}_k = \frac{\vec{u}_k}{\norm{\vec{u}_k}}
  \]}
  \onslide<+->{Resultado
  \(
    \{\vec{e}_1, \ldots, \vec{e}_n\}
  \)
  é uma \emph{base ortonormal}}
\end{frame}

%------------------------------------------------------------------------------%
\section{Exemplos}

%------------------------------------------------------------------------------%
\begin{question}
  Ortonormalize, pelo processo de Gram--Schmidt,
  \[
    \vec{v}_1 = (1,1,0)
    \qquad
    \vec{v}_2 = (1,0,1)
  \]
  e obtenha uma base ortonormal de
  \(
    \ger\{\vec{v}_1,\vec{v}_2\}
  \)
\end{question}

%------------------------------------------------------------------------------%
\begin{resolution}{Solução}
  \[
    \vec{u}_1
    \pause
    = \vec{v}_1
    \pause
    = (1, 1, 0)
  \]

  \pause
  \[
    \norm{\vec{u}_1}
    \pause
    = \sqrt{2}
  \]

  \pause
  \[
    \vec{e}_1
    \pause
    = \frac{1}{\sqrt{2}}(1,1,0)
  \]
\end{resolution}

%------------------------------------------------------------------------------%
\begin{resolution}{Solução}
\begin{columns}[t]
\column{0.3\textwidth}
  \begin{align*}
  \onslide<+->{\vec{u}_2}
  \onslide<+->{& = \vec{v}_2 - \frac{\vec{v}_2\cdot\vec{u}_1}{\norm{\vec{u}_1}^{2}} \; \vec{u}_1 \\}
  \onslide<+->{& = (1,0,1) - \frac{1}{2}(1,1,0) \\}
  \onslide<+->{& = \left( \frac{1}{2}, -\frac{1}{2}, 1 \right)}
  \end{align*}
\column{0.3\textwidth}
  \begin{align*}
    \onslide<+->{\norm{\vec{u}_2}^{2}}
    \onslide<+->{& = \frac{1}{4}+\frac{1}{4}+1\\}
    \onslide<+->{& = \frac{3}{2} \\}
    \onslide<+->{\norm{\vec{u}_2}}
    \onslide<+->{& = \frac{\sqrt{6}}{2}}
  \end{align*}
\column{0.3\textwidth}
  \begin{align*}
    \onslide<+->{\vec{e}_2}
    \onslide<+->{& = \frac{\vec{u}_2}{\norm{\vec{u}_2}}\\}
    \onslide<+->{& = \frac{1}{\sqrt{6}}(1,-1,2)}
  \end{align*}
\end{columns}
\end{resolution}

%------------------------------------------------------------------------------%
\begin{resolution}{Solução}
  \[
    \mathcal{B} =
    \left\{
      \frac{1}{\sqrt{2}}(1, 1, 0),
      \frac{1}{\sqrt{6}}(1, -1, 2)
    \right\}
  \]
\end{resolution}

%------------------------------------------------------------------------------%
\endinput

%------------------------------------------------------------------------------%
\begin{question}
  Ortonormalize a base
  \[
    B=
    \{
    (1,1,1),
    (1,0,1),
    (0,1,1)
    \}
  \]
  de $\R^3$ pelo processo de Gram--Schmidt
\end{question}

%------------------------------------------------------------------------------%
\begin{resolution}{Solução}
  \[
    \vec{u}_1
    \pause
    = (1, 1, 1)
  \]

  \pause
  \[
    \norm{\vec{u}_1}
    \pause
    = \sqrt{3}
  \]

  \pause
  \[
    \vec{e}_1
    \pause
    = \frac{1}{\sqrt{3}}(1,1,1)
  \]
\end{resolution}

%------------------------------------------------------------------------------%
\begin{resolution}{Solução}
\begin{columns}[t]
\column{0.4\textwidth}
  \begin{align*}
    \onslide<+->{\vec{u}_2}
    \onslide<+->{& = (1,0,1) - \frac{(1,0,1)\cdot(1,1,1)}{3}\, (1,1,1) \\}
    \onslide<+->{& = (1,0,1) - \frac{2}{3}\, (1,1,1) \\}
    \onslide<+->{& = \left(
        \frac{1}{3},
        -\frac{2}{3},
        \frac{1}{3}
      \right) \\}
    \onslide<+->{& = \frac{1}{3}(1,-2,1)}
  \end{align*}
\column{0.4\textwidth}
  \begin{align*}
    \onslide<+->{\norm{\vec{u}_2}}
    \onslide<+->{& = \frac{\sqrt{6}}{3}}
  \end{align*}
  \begin{align*}
    \onslide<+->{\vec{e}_2}
    \onslide<+->{& = \frac{\vec{u}_2}{\norm{\vec{u}_2}}\\}
    \onslide<+->{& = \frac{1}{\sqrt{6}}(1,-2,1)}
  \end{align*}
\end{columns}
\end{resolution}

%------------------------------------------------------------------------------%
\begin{resolution}{Solução}
\begin{align*}
  \onslide<+->{\vec{u}_3}
  \onslide<+->{& = (0,1,1)
    - \frac{2}{3}(1,1,1)
    -\frac{ (0,1,1)\cdot(1,-2,1) }{6} (1,-2,1)}
  \\
  \onslide<+->{& = (0,1,1)
    - \frac{2}{3}(1,1,1)
    + \frac{1}{6}(1,-2,1)}
  \\
  \onslide<+->{& = \left(
        -\frac{1}{2},
        0,
        \frac{1}{2}
      \right)}
  \\
  \onslide<+->{& = \frac{1}{2}(-1, 0, 1)}
\end{align*}
\end{resolution}

%------------------------------------------------------------------------------%
\begin{resolution}{Solução}
  \begin{align*}
    \onslide<+->{\norm{\vec{u}_3}}
    \onslide<+->{& = \frac{\sqrt{2}}{2}}
  \end{align*}
  \begin{align*}
    \onslide<+->{\vec{e}_3}
    \onslide<+->{& = \frac{\vec{u}_3}{\norm{\vec{u}_3}} \\}
    \onslide<+->{& = \frac{1}{\sqrt{2}}(-1, 0, 1)}
  \end{align*}

  \onslide<+->{\[
    \mathcal{B}
    = \left\{
      \frac{1}{\sqrt{3}}(1,1,1), \;
      \frac{1}{\sqrt{6}}(1,-2,1), \;
      \frac{1}{\sqrt{2}}(-1,0,1)
    \right\}
  \]}
\end{resolution}

%%------------------------------------------------------------------------------%
%\begin{resolution}{Verificação}
%\begin{align*}
%\vec{e}_1\cdot\vec{e}_2
%&=
%\frac{1-2+1}{\sqrt{18}}\\
%&=0\\
%\vec{e}_1\cdot\vec{e}_3
%&=
%\frac{-1+1}{\sqrt{6}}\\
%&=0\\
%\vec{e}_2\cdot\vec{e}_3
%&=
%\frac{-1+1}{\sqrt{12}}\\
%&=0
%\end{align*}
%
%\end{resolution}
%
%%------------------------------------------------------------------------------%
%\begin{resolution}{Verificação}
%\begin{align*}
%\norm{\vec{e}_1}^{2}
%&=
%\frac{3}{3}\\
%&=1\\
%\norm{\vec{e}_2}^{2}
%&=
%\frac{6}{6}\\
%&=1\\
%\norm{\vec{e}_3}^{2}
%&=
%\frac{2}{2}\\
%&=1
%\end{align*}
%\end{resolution}

%------------------------------------------------------------------------------%
\endinput

%%------------------------------------------------------------------------------%
\begin{question}
  Considere
  \[
  W=
  \left\{
  (x,y,z)\in\R^3
  :
  x+y+z=0
  \right\}
  \]
  Obtenha uma base ortonormal de $W$ pelo processo de Gram--Schmidt
\end{question}

%------------------------------------------------------------------------------%
\begin{resolution}{Solução}
Uma base de $W$ pode ser obtida de
\[
x=-y-z
\]
\[
(x,y,z)
=
y(-1,1,0)
+
z(-1,0,1)
\]
Escolha
\[
\vec{v}_1=(-1,1,0)
\qquad
\vec{v}_2=(-1,0,1)
\]
\end{resolution}

%------------------------------------------------------------------------------%
\begin{resolution}{Solução}
\begin{align*}
\vec{u}_1
&=(-1,1,0)\\
\norm{\vec{u}_1}
&=\sqrt{2}\\
\vec{e}_1
&=
\frac{1}{\sqrt{2}}(-1,1,0)
\end{align*}
\end{resolution}

%------------------------------------------------------------------------------%
\begin{resolution}{Solução}
\begin{align*}
\vec{u}_2
&=
(-1,0,1)
-
\frac{(-1,0,1)\cdot(-1,1,0)}
{2}
(-1,1,0)\\
&=
(-1,0,1)
-
\frac{1}{2}(-1,1,0)\\
&=
\left(
-\frac{1}{2},
-\frac{1}{2},
1
\right)\\
&=
\frac{1}{2}(-1,-1,2)
\end{align*}
\end{resolution}

%------------------------------------------------------------------------------%
\begin{resolution}{Solução}
\begin{align*}
\norm{\vec{u}_2}
&=
\frac{\sqrt{6}}{2}\\
\vec{e}_2
&=
\frac{1}{\sqrt{6}}(-1,-1,2)
\end{align*}
Portanto
\[
B_{\mathrm{on}}
=
\left\{
\frac{1}{\sqrt{2}}(-1,1,0)
,
\frac{1}{\sqrt{6}}(-1,-1,2)
\right\}
\]
\end{resolution}

%------------------------------------------------------------------------------%
\endinput

%------------------------------------------------------------------------------%
\begin{question}
  Determine as coordenadas de
  \[
    \vec{x} = (4, 2)
  \]
  na base ortonormal
  \[
    \mathcal{B} =
    \left\{
      \frac{1}{\sqrt{2}}(1,  1),
      \frac{1}{\sqrt{2}}(1, -1)
    \right\}
  \]
\end{question}

%------------------------------------------------------------------------------%
\begin{resolution}{Solução}
\begin{columns}[t]
\column{0.3\textwidth}
  \begin{align*}
    \onslide<+->{c_1}
    \onslide<+->{& = \vec{x}\cdot\vec{e}_1\\}
    \onslide<+->{& = (4, 2)\cdot\frac{1}{\sqrt{2}}(1,  1)\\}
    \onslide<+->{& = \frac{4+2}{\sqrt{2}}\\}
    \onslide<+->{& = \frac{2\times 3}{\sqrt{2}}\\}
    \onslide<+->{& = 3\sqrt{2}}
  \end{align*}
\column{0.3\textwidth}
  \begin{align*}
    \onslide<+->{c_2}
    \onslide<+->{& = \vec{x}\cdot\vec{e}_2\\}
    \onslide<+->{& = (4, 2)\cdot\frac{1}{\sqrt{2}}(1, -1)\\}
    \onslide<+->{& = \frac{4-2}{\sqrt{2}}\\}
    \onslide<+->{& = \frac{2}{\sqrt{2}}\\}
    \onslide<+->{& = \sqrt{2}}
  \end{align*}
\column{0.3\textwidth}
  \onslide<+->{\[
    [\vec{x}]_\mathcal{B}
    = \Vetor{3\sqrt{2}}{\sqrt{2}}
  \]}
\end{columns}
\end{resolution}

%%------------------------------------------------------------------------------%
%\begin{resolution}{Solução}
%\begin{align*}
%\norm{\vec{x}}^{2}
%&=
%4^{2}+2^{2}\\
%&=20\\
%c_1^{2}+c_2^{2}
%&=
%(3\sqrt{2})^{2}
%+
%(\sqrt{2})^{2}\\
%&=
%18+2\\
%&=20
%\end{align*}
%\end{resolution}

%------------------------------------------------------------------------------%
\endinput


%%------------------------------------------------------------------------------%
%\section{Lista Mínima}
%
%%------------------------------------------------------------------------------%
%\begin{frame}{Lista Mínima}
%  \begin{enumerate}
%    \item
%  \end{enumerate}
%
%  \vfill
%  Atenção: A prova é baseada no livro, não nas apresentações
%\end{frame}

%------------------------------------------------------------------------------%
\end{document}
%------------------------------------------------------------------------------%
